
\documentclass[aspectratio=169,xcolor=dvipsnames]{beamer}
\makeatletter
\def\input@path{{../../common/slides/theme/}{common/slides/theme/}{theme/}}
\makeatother
\usetheme{CleanEasy}
\usepackage[utf8]{inputenc}
\usepackage{lmodern}
\usepackage[T1]{fontenc}
\usepackage{fix-cm}
\usepackage{amsmath}
\usepackage{mathtools}
\usepackage{listings}
\usepackage{xcolor}
\usepackage{hyperref}
\usepackage{graphicx}
\graphicspath{{../../common/slides/assets/}{common/slides/assets/}{./}}
\usepackage{booktabs}
\usepackage{tikz}
\usetikzlibrary{positioning, shapes, arrows, calc, decorations.pathreplacing, arrows.meta, backgrounds, patterns, overlay-beamer-styles, angles, quotes, intersections}
\usepackage{etoolbox}
\usepackage{animate}

\lstset{
  basicstyle=\ttfamily\small,
  keywordstyle=\color{blue},
  commentstyle=\color{green!60!black},
  stringstyle=\color{red},
  showstringspaces=false,
  breaklines=true,
  frame=single,
  rulecolor=\color{black!30},
  backgroundcolor=\color{black!5},
  numbers=left,
  numberstyle=\tiny\color{black!70},
  numbersep=5pt
}
% big first triple, smaller multiples underneath
\newcommand{\triplestack}[3]{%
\begin{tabular}{@{}c@{}}
{\Large #1}\\[0.2em]
{\scriptsize #2}\\[-0.2em]
{\scriptsize #3}\\
\vdots
\end{tabular}%
}



%----------------------------------------------------------------------------------------
%    TITLE PAGE
%----------------------------------------------------------------------------------------


% why were these two topics combined? no one will know.
\title[Quads \& some Trig]{Quads \& some Trig}

\author[IA Math Team]{IA Math Team}

\vspace{-2cm}\date{Feb. 12, 2026}
%---------------------------------------------------------------------------------------

\begin{document}

\begin{frame}[t]
  \titlepage
\end{frame}

% FLOYD COUNTY ANNOUNCEMENTS!!!
\section{LISTEN UP!}

\section{Warmup...}


\begin{frame}[t]{2022 AMC 10A, \# 15}
Quadrilateral $ABCD$ with side lengths $AB=7, BC=24, CD=20, DA=15$ is inscribed in a circle. The area interior to the circle but exterior to the quadrilateral can be written in the form $\frac{a\pi-b}{c},$ where $a,b,$ and $c$ are positive integers such that $a$ and $c$ have no common prime factor. What is $a+b+c?$
\end{frame}


\begin{frame}[t]{Pythagorean Triples Time!!!}
How many do you know?
\end{frame}

\begin{frame}[t]{Pythagorean Triples Time!!!}
\centering
\scriptsize
\setlength{\tabcolsep}{15pt}
\renewcommand{\arraystretch}{1.5}

\begin{tabular}{cccc}
\triplestack{$(3,4,5)$}{$(6,8,10)$}{$(9,12,15)$} &
\triplestack{$(5,12,13)$}{$(10,24,26)$}{$(15,36,39)$} &
\triplestack{$(8,15,17)$}{$(16,30,34)$}{$(24,45,51)$} &
\triplestack{$(7,24,25)$}{$(14,48,50)$}{$(21,72,75)$}
\end{tabular}

\vspace{2.6em}

\begin{tabular}{cccc}
\triplestack{$(20,21,29)$}{$(40,42,58)$}{$(60,63,87)$} &
\triplestack{$(12,35,37)$}{$(24,70,74)$}{$(36,105,111)$} &
\triplestack{$(9,40,41)$}{$(18,80,82)$}{$(27,120,123)$} &
\triplestack{$(28,45,53)$}{$(56,90,106)$}{$(84,135,159)$}
\end{tabular}

\end{frame}

\begin{frame}[t]{Pythagorean Triples Time!!!}

\begin{itemize}
    \item It's not too important to know \textit{a ton} of triples (as opposed to your powers): the real test of skill becomes how you can:
    \begin{itemize}
        \item \textbf{\textit{Identify them in the wild}} (just like we did)
        \begin{itemize}
            \item Use them to show that some triangle is right
            \item Use two to complete a triple
        \end{itemize}
    \end{itemize}
\end{itemize}

% Anyways, enough for triples and back to the actual topic of the day, which is quadrilaterals. 

% Does anyone wanna ask why? The reason is because if I made the topic on Pythagorean triples, you would immediately know to use them in otherwise tricky problems - does that make sense? You guys will just be working on a worksheet on quadrilaterals this practice.

% Some of these are advanced, such as knowing some concepts. We'll go around the classroom to gather small groups working on the same topic. 

\end{frame}






\end{document}